\documentclass[11pt]{article}
\usepackage[margin=1in]{geometry}
\usepackage{amsmath,amssymb,amsthm}
\usepackage[htt]{hyphenat}
\usepackage[hidelinks]{hyperref}
\usepackage{booktabs}
\usepackage{listings}
\usepackage{xcolor}
\hyphenation{ap-prox-i-mant telescop-ing can-cel-la-tion}
\sloppy

\newtheorem{theorem}{Theorem}[section]
\newtheorem{lemma}[theorem]{Lemma}
\newtheorem{proposition}[theorem]{Proposition}
\newtheorem{corollary}[theorem]{Corollary}
\theoremstyle{definition}
\newtheorem{definition}[theorem]{Definition}
\newtheorem{remark}[theorem]{Remark}

\newcommand{\E}{\mathbb{E}}
\newcommand{\R}{\mathbb{R}}
\newcommand{\dd}{\,\mathrm{d}}
\DeclareMathOperator{\cv}{\leq_{\mathrm{cx}}}

\title{Higher-Order Cancellation in Exponential Approximants:\\
An Explicit $O(n^{-4})$ Telescoping Family}
\author{Josh Bald\\ \small Independent Researcher\\ \small jpbald93@gmail.com}
\date{September 2026 \quad (V5)}

\begin{document}
\maketitle

\begin{abstract}
We study elementary sequences $(x_n)$ converging to exponential constants $e^{c}$
for which the successive differences $D_n := x_{n+1}-x_n$ form a rapidly decaying
telescoping tail. For any convergent sequence $(x_n)\to L$, the identity
$L = x_1 + \sum_{n\ge1} D_n$ reduces quantitative convergence questions to
asymptotics of $D_n$; in particular, if $D_n = K n^{-(k+1)} + O(n^{-(k+2)})$ then
$L - x_N = \tfrac{K}{k} N^{-k} + O(N^{-(k+1)})$. Within the constrained ansatz
\[
  x_n = \Bigl(1 + \frac{a_1}{n} + \frac{a_2}{n^2}\Bigr)^{cn+\gamma},
\]
we use logarithmic expansions and coefficient cancellation to enforce higher-order
decay of $D_n$. The methodological ingredients are classical; the point here is
that, for this specific quadratic-base/affine-exponent competitor class, the
order-$3$ constraints admit an explicit closed-form solution.

Our main result is a two-branch family achieving telescoping order $3$
(i.e.\ $|D_n| = \Theta(n^{-4}))$ for any real $c \neq 0$:
\[
  x_n^{(\mathrm{I})} = \Bigl(1+\frac1n+\frac{\sqrt3}{6n^2}\Bigr)^{cn+c(\frac12-\frac{\sqrt3}{6})},
  \qquad
  x_n^{(\mathrm{II})} = \Bigl(1+\frac1n-\frac{\sqrt3}{6n^2}\Bigr)^{cn+c(\frac12+\frac{\sqrt3}{6})}.
\]
We prove $\lim_{n\to\infty} x_n^{(\cdot)} = e^{c}$ and derive explicit signed
asymptotics
\[
  D_n^{(\cdot)} \sim -\,\frac{3e^{c}A_3^{(\cdot)}}{n^{4}},
  \qquad
  A_3^{(\mathrm{I})} = \frac{c\,(2\sqrt3-3)}{72},\quad
  A_3^{(\mathrm{II})} = -\frac{c\,(3+2\sqrt3)}{72},
\]
implying the tail error $|x_N - e^{c}| = \Theta(N^{-3})$ for $c\neq0$. (At $c=0$
both sequences are identically $1$ and the asymptotic statements are vacuous.)
\end{abstract}

\section{Introduction}

\subsection{Telescoping representations}
For any convergent sequence $(x_n)\to L$, we have the telescoping identity
\[
  L = x_1 + \sum_{n=1}^{\infty}(x_{n+1}-x_n) = x_1 + \sum_{n=1}^{\infty} D_n .
\]
The speed of convergence is governed by the decay of the differences $D_n$;
faster decay of $D_n$ yields a smaller tail
\[
  R_N := \sum_{n>N} D_n = L - x_{N+1}.
\]

\subsection{Positivity of the bases (well-defined real powers)}
Because we use real (not complex) exponentiation, we record a positivity check
ensuring the bases remain positive for all $n\ge1$.

\begin{lemma}[Positivity of the bases]\label{lem:pos}
For all integers $n\ge1$,
\[
  1+\frac1n+\frac{\sqrt3}{6n^2} > 0
  \quad\text{and}\quad
  1+\frac1n-\frac{\sqrt3}{6n^2} > 0 .
\]
Hence $x_n^{(\mathrm{I})}$ and $x_n^{(\mathrm{II})}$ are well-defined as real
numbers for all $n\ge1$.
\end{lemma}

\begin{proof}
The first inequality is immediate since $1+\tfrac1n+\tfrac{\sqrt3}{6n^2}>1$.
For the second, for $n\ge1$ we have $n^{-2}\le1$, hence
$1+\tfrac1n-\tfrac{\sqrt3}{6n^2}\ge 1-\tfrac{\sqrt3}{6}=\tfrac{6-\sqrt3}{6}>0$.
\end{proof}

\subsection{From difference decay to error decay}
The following standard comparison makes the meaning of ``telescoping order''
explicit.

\begin{lemma}[Tail rate from difference rate]\label{lem:tail}
Let $(x_n)$ converge to $L$ and set $D_n := x_{n+1}-x_n$. Suppose that for some
$k\ge1$,
\[
  D_n = K n^{-(k+1)} + O(n^{-(k+2)}) \qquad (K\neq0).
\]
Then
\[
  L - x_N = \sum_{n\ge N} D_n = \frac{K}{k}\,N^{-k} + O(N^{-(k+1)}).
\]
In particular, telescoping order $k$ (i.e.\ $D_n=\Theta(n^{-(k+1)})$) corresponds
to an $O(N^{-k})$ convergence rate for the error $|x_N-L|$.
\end{lemma}

\begin{proof}
This follows from the telescoping identity and the asymptotic evaluation of the
tail of a $p$-series,
$\sum_{n\ge N} n^{-(k+1)} = \tfrac1k N^{-k} + O(N^{-(k+1)})$.
\end{proof}

\begin{remark}[Signed asymptotics vs.\ absolute bounds]\label{rem:theta}
Lemma~\ref{lem:tail} is stated with a \emph{signed} asymptotic for $D_n$; the
conclusion $L-x_N=\Theta(N^{-k})$ uses that sign. A bare absolute bound
$|D_n|=\Theta(n^{-(k+1)})$ does \emph{not} by itself force a $\Theta$ tail rate:
e.g.\ $x_n = 1+(-1)^n n^{-4}$ has $|D_n|=\Theta(n^{-4})$ but $|x_n-1|=n^{-4}$,
not $\Theta(n^{-3})$. Throughout this paper the relevant families admit genuine
signed power-series expansions, so Lemma~\ref{lem:tail} applies directly.
\end{remark}

\subsection{Why telescoping matters (context and applications)}
Although the present paper focuses on elementary closed forms for $e^{c}$, the
telescoping viewpoint is broadly useful as a way to translate convergence into an
explicit remainder model. Many classic acceleration techniques can be viewed as
constructing and cancelling leading terms in an error expansion, and the NIST DLMF
provides standard language and examples for sequence transformations that
accelerate convergence (Aitken, Shanks/Wynn-$\varepsilon$, Levin--Weniger,
etc.)~\cite{dlmf39} (see also~\cite{corless2023}). In a more recent probabilistic
framing, Oates--Karvonen et al.\ propose probabilistic Richardson extrapolation,
where an error model is fitted/propagated to accelerate convergence while
quantifying uncertainty~\cite{richardson2024,richardson2024arxiv}. In this light,
explicit telescoping families with known $D_n$ asymptotics provide a clean
laboratory for validating order-detection by regression and illustrating how
algebraic constraints enforce higher-order cancellation in a controlled setting.

\subsection{Classical examples and the classical design principle}
The prototypical example is Euler's sequence $(1+1/n)^n\to e$. More generally,
for $(1+a/n)^{bn}\to e^{ab}$, a direct expansion of the logarithm
(Proposition~\ref{prop:order1} below) gives
\begin{equation}\label{eq:classical}
  D_n = \frac{e^{ab}\,a^2 b}{2n^2} + O(n^{-3}),
\end{equation}
so the leading term vanishes \emph{iff} $a=0$ or $b=0$ (both of which give the
constant sequence $1$). In particular $(1+c/n)^{cn}$ has limit $e^{c^2}$ and
difference order exactly $\Theta(n^{-2})$ for $c\neq0$; at $c=1$ it is Euler's
sequence. (An earlier draft of this work asserted the stronger and incorrect
claim that the leading term cancels at $a=b$; see
Remark~\ref{rem:correction}.) The methodology used throughout this
paper---taking logarithms, expanding in powers of $1/n$, and enforcing coefficient
cancellations to increase decay order---is standard asymptotic-expansion practice;
see, for example,~\cite{flajolet}.

\subsection{Competitor class and novelty boundary}
\begin{definition}[Competitor class: quadratic base, affine exponent]\label{def:comp}
We call a sequence \emph{quadratic-base, affine-exponent} if it has the form
\[
  x_n = \Bigl(1 + \frac{a_1}{n} + \frac{a_2}{n^2}\Bigr)^{cn+\gamma},
\]
with parameters $a_1,a_2,\gamma$ given in elementary closed form (and $c$ a
prescribed real constant).
\end{definition}

Within this competitor class, telescoping-order constraints reduce to a finite
algebraic system in $(a_1,a_2,\gamma)$ (after the normalizations described later).
The novelty boundary in this paper is therefore not the existence of coefficient
matching, but the fact that in the quadratic-base, affine-exponent class one can
solve the order-$3$ constraints explicitly in closed form and obtain a two-branch
parameterization together with explicit leading asymptotic constants.

\subsection{Relation to prior work and novelty (conservative framing)}
Refined investigations of exponential constants and limiting processes (including
$e$ and the Euler--Mascheroni constant $\gamma$) have a long history; see
Glaisher~\cite{glaisher}, Adams~\cite{adams}, and Mitchell--Strain~\cite{mitchell}.
At a methodological level, logarithmic expansions and cancellation arguments are
classical; see Brent~\cite{brent} and~\cite{flajolet}. Modern
convergence-acceleration context can be found in the DLMF
overview~\cite{dlmf39} and Corless~\cite{corless2023}. The present work also builds
on the telescoping-difference framework developed in~\cite{balditerative}.

Accordingly, the novelty of this paper does not lie in asymptotic expansions per
se, but in producing a fully explicit closed-form order-$3$ telescoping family
inside the tightly constrained competitor class of Definition~\ref{def:comp},
including explicit signed leading coefficients controlling $D_n$ and a direct
implication for the $N^{-3}$ error rate.

\subsection{Main result}
\begin{theorem}[Explicit order-$3$ family in a quadratic-base, affine-exponent ansatz]\label{thm:main}
Fix $c\in\R\setminus\{0\}$. Define
\[
  x_n^{(\mathrm{I})} := \Bigl(1+\frac1n+\frac{\sqrt3}{6n^2}\Bigr)^{cn+c(\frac12-\frac{\sqrt3}{6})},
  \qquad
  x_n^{(\mathrm{II})} := \Bigl(1+\frac1n-\frac{\sqrt3}{6n^2}\Bigr)^{cn+c(\frac12+\frac{\sqrt3}{6})}.
\]
Then for each family:
\begin{enumerate}
  \item $\displaystyle \lim_{n\to\infty} x_n^{(\cdot)} = e^{c}$.
  \item $D_n^{(\cdot)} := x_{n+1}^{(\cdot)} - x_n^{(\cdot)}$ satisfies
        $|D_n^{(\cdot)}| = \Theta(n^{-4})$.
\end{enumerate}
More precisely,
\begin{equation}\label{eq:mainasy}
  D_n^{(\cdot)} \sim -\,\frac{3e^{c}A_3^{(\cdot)}}{n^{4}}
  \qquad\text{as } n\to\infty,
\end{equation}
where
\[
  A_3^{(\mathrm{I})} = \frac{c\,(2\sqrt3-3)}{72},
  \qquad
  A_3^{(\mathrm{II})} = -\frac{c\,(3+2\sqrt3)}{72},
\]
so $K^{(\cdot)} = 3e^{c}\,|A_3^{(\cdot)}|$.
\end{theorem}

\begin{remark}[Branch labeling and sign]\label{rem:branch}
We label branches by the magnitude of the leading asymptotic coefficient:
$x_n^{(\mathrm{I})}$ has smaller $|D_n|$ (faster convergence) and
$x_n^{(\mathrm{II})}$ has larger $|D_n|$ (slower convergence), since
$0<2\sqrt3-3<3+2\sqrt3$. For $c>0$ the two leading errors have \emph{opposite
signs}: $x_n^{(\mathrm{I})}$ approaches $e^{c}$ from above (decreasing), while
$x_n^{(\mathrm{II})}$ approaches $e^{c}$ from below (increasing); the signs
reverse for $c<0$.
\end{remark}

\section{Framework and Asymptotic Expansion}

\subsection{General setup}
We consider sequences of the form
\[
  x_n = f(1/n)^{g(n)}, \qquad
  f(t) = 1 + \alpha_1 t + \alpha_2 t^2 + \cdots, \qquad
  g(n) = \beta n + \gamma .
\]

\begin{definition}[Telescoping order]\label{def:order}
A sequence $(x_n)\to L$ has \emph{telescoping order} $k$ if
$|D_n| = \Theta(n^{-(k+1)})$.
\end{definition}

\begin{remark}\label{rem:order}
The phrase ``order $k$'' here refers to the decay rate
$D_n\asymp n^{-(k+1)}$. By Lemma~\ref{lem:tail}, this corresponds to an
$O(N^{-k})$ rate for the approximation error $|x_N-L|$ (with a signed
asymptotic, as discussed in Remark~\ref{rem:theta}).
\end{remark}

\subsection{Asymptotic expansion of $\log x_n$}
Write $u_n := f(1/n)-1 = \alpha_1/n + \alpha_2/n^2 + \cdots$. Then
$\log(1+u_n) = u_n - \tfrac{u_n^2}{2} + \tfrac{u_n^3}{3} - \cdots$ and
$\log x_n = g(n)\log(1+u_n)$. Expanding in powers of $1/n$,
\begin{equation}\label{eq:logexp}
  \log x_n = c_0 + \frac{A_1}{n} + \frac{A_2}{n^2} + \frac{A_3}{n^3} + \frac{A_4}{n^4} + O(n^{-5}),
\end{equation}
where $c_0,A_1,A_2,\dots$ are explicit polynomials in $\{\alpha_j,\beta,\gamma\}$
(given in \S\ref{sec:order3}).

\subsection{From $\log x_n$ to $D_n$}
The following lemma is the bridge from the logarithmic expansion to the
difference asymptotic. Note the additional regularity hypothesis: a bare
$O(n^{-(k+2)})$ remainder cannot be differenced as though it were smooth (the
remainder's finite difference need not be $O(n^{-(k+3)})$).

\begin{lemma}[Log-expansion implies difference asymptotic]\label{lem:logdiff}
Suppose that for some integer $k\ge0$ the sequence admits an asymptotic
expansion in powers of $1/n$,
\begin{equation}\label{eq:hyp}
  \log x_n = c + \frac{A}{n^{k+1}} + \frac{B}{n^{k+2}} + O(n^{-(k+3)}),
  \qquad A\neq0 .
\end{equation}
Then
\[
  x_n = e^{c}\Bigl(1 + \frac{A}{n^{k+1}} + O(n^{-(k+2)})\Bigr),
  \qquad
  D_n := x_{n+1}-x_n = -(k+1)e^{c}\,A\,n^{-(k+2)} + O(n^{-(k+3)}).
\]
\end{lemma}

\begin{proof}
Exponentiating~\eqref{eq:hyp} and using $\exp(z)=1+z+O(z^2)$ gives
$x_n = e^{c}\bigl(1 + A n^{-(k+1)} + B n^{-(k+2)} + O(n^{-(k+3)})\bigr)$ for
$k\ge1$, since $(n^{-(k+1)})^2 = O(n^{-(k+3)})$; for $k=0$ the $n^{-2}$
coefficient instead becomes $B+A^2/2$, which is immaterial at the order
computed. Writing $x_n = e^{c}\bigl(1 + A n^{-(k+1)} + R_n\bigr)$ we thus have
$R_n = B n^{-(k+2)} + O(n^{-(k+3)})$, so
\[
  \Delta R_n = B\bigl[(n+1)^{-(k+2)} - n^{-(k+2)}\bigr] + O(n^{-(k+3)})
            = O(n^{-(k+3)})
\]
(the mean value theorem gives the first term, and the triangle inequality
bounds the remainder). Hence
\[
  D_n = e^{c}A\bigl[(n+1)^{-(k+1)} - n^{-(k+1)}\bigr] + O(n^{-(k+3)}),
\]
and $(n+1)^{-(k+1)} - n^{-(k+1)} = -(k+1)n^{-(k+2)} + O(n^{-(k+3)})$ by the mean
value theorem, which completes the proof.
\end{proof}

\emph{Design principle.} To achieve telescoping order $k$ in the sense of
Definition~\ref{def:order}, enforce $c_0=c$ and $A_1=\cdots=A_{k-1}=0$, leaving
$A_k\neq0$. (Indeed, if the first nonzero coefficient is $A_k$, then
$\log x_n = c + A_k n^{-k} + O(n^{-(k+1)})$, and
Lemma~\ref{lem:logdiff} with the relabeled index $k-1$ gives
$D_n = -k e^{c}A_k n^{-(k+1)} + O(n^{-(k+2)}) = \Theta(n^{-(k+1)})$, i.e.\ order
$k$.)

\section{Order-1: Classical Cancellation}
\begin{proposition}[Two-parameter family]\label{prop:order1}
For $x_n = (1+a/n)^{bn}$ with real $a,b$ and $n$ sufficiently large that the base
is positive,
\[
  \lim_{n\to\infty} x_n = e^{ab},
  \qquad
  D_n = \frac{e^{ab}\,a^2 b}{2n^2} + O(n^{-3}).
\]
Consequently $D_n = O(n^{-3})$ if and only if $a=0$ or $b=0$ (in which case
$x_n\equiv1$); for $a,b\neq0$ the difference order is exactly $\Theta(n^{-2})$.
\end{proposition}

\begin{proof}
Since $\log(1+a/n) = a/n - a^2/(2n^2) + O(n^{-3})$,
\[
  \log x_n = bn\log\!\Bigl(1+\frac{a}{n}\Bigr)
           = ab - \frac{a^2 b}{2n} + \frac{a^3 b}{3n^2} + O(n^{-3}).
\]
Hence $x_n = e^{ab}\bigl(1 - \tfrac{a^2 b}{2n} + O(n^{-2})\bigr)$, and
Lemma~\ref{lem:logdiff} with $k=0$ and $A=-a^2b/2$ gives
$D_n = e^{ab}\,\tfrac{a^2 b}{2n^2} + O(n^{-3})$.
\end{proof}

\begin{remark}[Correction of an earlier draft]\label{rem:correction}
An earlier version of this paper stated
$D_n = e^{ab}ab(a-b)/(2n^2)+O(n^{-3})$ and concluded that the leading term
cancels at $a=b$. That is incorrect: the leading coefficient is $a^2b/2$, not
$ab(a-b)/2$, and it does not vanish at $a=b$. Numerically, for Euler's sequence
($a=b=1$) one has $n^2 D_n/e \to 1/2$, i.e.\ $D_n=\Theta(n^{-2})$, not
$O(n^{-3})$. The corrected statement is Proposition~\ref{prop:order1}; see also
Table~\ref{tab:compare}.
\end{remark}

\section{Explicit Order-3 Constructions}\label{sec:order3}

\subsection{Quadratic base, affine exponent}
Consider
\[
  x_n = \Bigl(1 + \frac{\alpha_1}{n} + \frac{\alpha_2}{n^2}\Bigr)^{\beta n+\gamma}.
\]
Let $u_n = \alpha_1/n + \alpha_2/n^2$. Expanding $\log(1+u_n)$ and multiplying by
$\beta n + \gamma$ yields
\begin{align}
  c_0 &= \alpha_1\beta, \label{eq:c0}\\
  A_1 &= \alpha_1\gamma + \beta\alpha_2 - \frac{\alpha_1^2\beta}{2}, \label{eq:A1}\\
  A_2 &= \gamma\alpha_2 - \frac{\alpha_1^2\gamma}{2} + \frac{\alpha_1^3\beta}{3} - \alpha_1\alpha_2\beta, \label{eq:A2}\\
  A_3 &= \gamma\Bigl(\frac{\alpha_1^3}{3} - \alpha_1\alpha_2\Bigr)
        + \beta\Bigl(\alpha_1^2\alpha_2 - \frac{\alpha_1^4}{4} - \frac{\alpha_2^2}{2}\Bigr), \label{eq:A3}\\
  A_4 &= \gamma\Bigl(\alpha_1^2\alpha_2 - \frac{\alpha_1^4}{4} - \frac{\alpha_2^2}{2}\Bigr)
        + \beta\Bigl(\frac{\alpha_1^5}{5} - \alpha_1^3\alpha_2 + \alpha_1\alpha_2^2\Bigr). \label{eq:A4}
\end{align}
(The $A_4/n^4$ term is quoted for completeness; it controls the subleading
correction and is used in \S\ref{sec:sub}.)

\begin{remark}[Correction of an earlier draft]\label{rem:eq4}
An earlier version printed the $\beta$-part of~\eqref{eq:A3} as
$\beta(\alpha_1^4/4 - \alpha_1^2\alpha_2/2)$, with two incorrect signs and the
$- \alpha_2^2/2$ term missing. The correct form is~\eqref{eq:A3}. As a check,
substituting the branch parameters into the incorrect form gives
$c(9-4\sqrt3)/18$, which agrees neither with the correct value nor with the
earlier draft's own stated constant.
\end{remark}

\subsection{Constraint system and closed-form solutions}
To achieve telescoping order $3$ we require $c_0=c$, $A_1=0$, $A_2=0$. We have
four parameters $\{\alpha_1,\alpha_2,\beta,\gamma\}$ and three constraints. Fix
$\alpha_1=1$ and $\beta=c$ to satisfy $c_0=c$. The remaining constraints yield
\[
  \gamma + c\alpha_2 - \frac{c}{2} = 0,
  \qquad
  \gamma\alpha_2 - \frac{\gamma}{2} + \frac{c}{3} - c\alpha_2 = 0.
\]
From the first, $\gamma = c(\tfrac12 - \alpha_2)$. Substituting into the second
gives $\alpha_2^2 = \tfrac1{12}$, so $\alpha_2 = \pm\tfrac{\sqrt3}{6}$. Thus we
obtain two explicit branches:
\[
  \text{Family I: } \alpha_2 = +\frac{\sqrt3}{6},\quad
    \gamma = c\Bigl(\frac12 - \frac{\sqrt3}{6}\Bigr);
  \qquad
  \text{Family II: } \alpha_2 = -\frac{\sqrt3}{6},\quad
    \gamma = c\Bigl(\frac12 + \frac{\sqrt3}{6}\Bigr).
\]

\subsection{Explicit formula for $A_3$ and nonvanishing}\label{sec:sub}
Substituting $\alpha_1=1$, $\beta=c$ and the corresponding $\gamma$ values
into~\eqref{eq:A3} gives
\[
  A_3^{(\mathrm{I})} = \frac{c\,(2\sqrt3-3)}{72},
  \qquad
  A_3^{(\mathrm{II})} = -\frac{c\,(3+2\sqrt3)}{72}.
\]
Since $2\sqrt3-3>0$ and $3+2\sqrt3>0$, we have $A_3^{(\cdot)}\neq0$ whenever
$c\neq0$ (with opposite signs on the two branches). By
Lemma~\ref{lem:logdiff} (with $k=2$ and $A=A_3$),
\[
  D_n \sim -\,\frac{3e^{c}A_3}{n^{4}} .
\]
For completeness, substituting into~\eqref{eq:A4} gives the subleading
coefficients
\[
  A_4^{(\mathrm{I})} = \frac{c\,(39-25\sqrt3)}{720},
  \qquad
  A_4^{(\mathrm{II})} = \frac{c\,(39+25\sqrt3)}{720},
\]
which quantify the $O(n^{-5})$ corrections to $D_n$ and address
Open Problem~1 below.

\begin{remark}[Correction of an earlier draft]\label{rem:const}
An earlier version stated $A_3^{(\mathrm{I})}=c(9-5\sqrt3)/36$ and
$A_3^{(\mathrm{II})}=c(9+5\sqrt3)/36$, together with the multiplier $-4e^{c}$
in~\eqref{eq:mainasy}. Both are incorrect: the constants are as above, and the
multiplier is $-3e^{c}$ (as required by Lemma~\ref{lem:logdiff} with $k=2$). The
paper's own numerical tables already reflected the corrected values; only the
closed forms were in error.
\end{remark}

\subsection{Proof of the main theorem}
\begin{proof}[Proof of Theorem~\ref{thm:main}]
With the parameters above, we have $c_0=c$, $A_1=0$, $A_2=0$ and the stated
$A_3\neq0$. Hence
\[
  \log x_n = c + \frac{A_3}{n^3} + \frac{A_4}{n^4} + O(n^{-5}),
\]
which is of the form~\eqref{eq:hyp} with $k=2$, $A=A_3$, $B=A_4$.
Lemma~\ref{lem:logdiff} yields $D_n \sim -3e^{c}A_3 n^{-4}$, so
$|D_n|=\Theta(n^{-4})$. Finally, Lemma~\ref{lem:tail} implies the corresponding
error rate $|x_N-e^{c}| = \Theta(N^{-3})$.
\end{proof}

\section{Numerical Validation}
We implement $x_n^{(\cdot)}$ at $150$-digit precision using \texttt{mpmath} and
compute $D_n = x_{n+1}-x_n$ for $100\le n\le 10000$. Log--log regression over the
grid $n=100,125,\dots,10000$ yields
\[
  \text{Family I: } \log|D_n| = -3.9973\,\log n - 2.9696,\quad R^2=0.9999997,
\]
\[
  \text{Family II: } \log|D_n| = -3.9969\,\log n - 0.3391,\quad R^2=0.9999996,
\]
consistent with the predicted $n^{-4}$ decay. (The fitted slopes differ from
$-4$ by a small finite-range curvature; the intercepts likewise carry an
$O(n^{-1})$ correction. The numbers above are regenerated by the script in
Appendix~\ref{app:code}.)

\subsection{A log--log ``figure'' (self-contained, text-only)}
A literal log--log plot of $|D_n|$ versus $n$ would appear as an approximately
straight line with slope near $-4$ for both branches. The regression above
provides a compact self-contained surrogate for the visual.

Detailed convergence for $c=1$ (target $e\approx2.71828$):

\begin{table}[h]\centering
\begin{tabular}{rrrrr}
\toprule
$n$ & $|x_n^{(\mathrm{I})}-e|$ & $|D_n^{(\mathrm{I})}|$ & $|x_n^{(\mathrm{II})}-e|$ & $|D_n^{(\mathrm{II})}|$\\
\midrule
100  & $1.74\times10^{-8}$  & $5.09\times10^{-10}$ & $2.41\times10^{-7}$  & $7.06\times10^{-9}$\\
500  & $1.40\times10^{-10}$ & $8.36\times10^{-13}$ & $1.95\times10^{-9}$  & $1.16\times10^{-11}$\\
1000 & $1.75\times10^{-11}$ & $5.24\times10^{-14}$ & $2.44\times10^{-10}$ & $7.29\times10^{-13}$\\
5000 & $1.40\times10^{-13}$ & $8.40\times10^{-17}$ & $1.95\times10^{-12}$ & $1.17\times10^{-15}$\\
\bottomrule
\end{tabular}
\caption{Convergence of $x_n^{(\cdot)}$ to $e$ ($c=1$) at 150-digit precision.}
\label{tab:c1}
\end{table}

\subsection{Generality check: additional values of $c$}
Table~\ref{tab:gen} reports representative errors and differences for $c=2$ and
$c=-1$.

\begin{table}[h]\centering
\begin{tabular}{rrrrr}
\toprule
$c$ & $n$ & $|x_n^{(\mathrm{I})}-e^{c}|$ & $|D_n^{(\mathrm{I})}|$ & $|D_n^{(\mathrm{II})}|$\\
\midrule
$2$  & 1000 & $9.52\times10^{-11}$ & $2.85\times10^{-13}$ & $3.97\times10^{-12}$\\
$2$  & 5000 & $7.62\times10^{-13}$ & $4.57\times10^{-16}$ & $6.36\times10^{-15}$\\
$-1$ & 1000 & $2.37\times10^{-12}$ & $7.09\times10^{-15}$ & $9.87\times10^{-14}$\\
$-1$ & 5000 & $1.90\times10^{-14}$ & $1.14\times10^{-17}$ & $1.58\times10^{-16}$\\
\bottomrule
\end{tabular}
\caption{Illustrative generality check for $c=2$ and $c=-1$ at high precision.}
\label{tab:gen}
\end{table}

\section{Comparison with Existing Methods}

\subsection{Convergence-rate comparison}
By Lemma~\ref{lem:tail}, if $D_n$ has the signed asymptotic
$D_n = K n^{-(k+1)} + O(n^{-(k+2)})$ with $K\neq0$ (as holds for every family
compared below), then $|x_N-L|=\Theta(N^{-k})$. Thus, for a target error
$\varepsilon$, one expects $N\asymp\varepsilon^{-1/k}$ in such telescoping
constructions.

\begin{table}[h]\centering
\begin{tabular}{llll}
\toprule
Method & Difference decay $D_n$ & Error decay $|x_N-L|$ & $N$ for $|x_N-L|<\varepsilon$\\
\midrule
Euler $(1+1/n)^n$             & $\Theta(n^{-2})$ & $\Theta(N^{-1})$ & $\asymp\varepsilon^{-1}$\\
Classical $(1+c/n)^{cn}$      & $\Theta(n^{-2})$ & $\Theta(N^{-1})$ & $\asymp\varepsilon^{-1}$\\
Order-3 (this work)           & $\Theta(n^{-4})$ & $\Theta(N^{-3})$ & $\asymp\varepsilon^{-1/3}$\\
\bottomrule
\end{tabular}
\caption{Telescoping decay of differences and implied error rates (parameterized rows: fixed $c\neq0$).}
\label{tab:compare}
\end{table}

\begin{remark}\label{rem:table3}
An earlier draft listed the middle row as $\Theta(n^{-3})/\Theta(N^{-2})$ with
target $e^{c}$; that is incorrect. For $(1+c/n)^{cn}$ the limit is $e^{c^2}$ and
the difference order is $\Theta(n^{-2})$ (Proposition~\ref{prop:order1}); at
$c=1$ it coincides with Euler's sequence in the first row. To obtain
$O(N^{-2})$ error decay one must instead \emph{enlarge} the family (e.g.\ by an
affine exponent); by Proposition~\ref{prop:order1} no nonconstant member of
$(1+a/n)^{bn}$ has order $2$, so this is not pursued here.
\end{remark}

\subsection{Computational cost and limitations (pedagogical framing)}
We make no competitive complexity claim. For a fixed working precision the
construction attains tail error $<\varepsilon$ at index $N\asymp\varepsilon^{-1/3}$,
so a single evaluation of $x_N$ suffices; the telescoping identity means one
need not sum all $N$ differences. We deliberately state no bit-complexity bound:
the cost of a real power $a^{b}$ at precision $p$ depends on the underlying
$\exp/\log$ algorithm and is not $O(\log n)$ arithmetic operations in general. This is not
intended as a competitive high-precision algorithm for computing $e^{c}$ (where
AGM-based and related methods converge far faster). Instead, the value of the
construction is primarily structural and pedagogical: it provides a compact
closed-form example where higher-order cancellation can be derived, proved, and
numerically verified end-to-end using only elementary expansions and a small
parameter system. In particular, it is well-suited for illustrating (i) how
asymptotic expansions control discrete differences, (ii) how cancellation
constraints translate into algebraic equations, and (iii) how regression on
log--log data can detect decay exponents in practice.

\section{Open Problems}
\begin{enumerate}
  \item \textbf{Subleading corrections:} the $O(n^{-5})$ and higher terms are now
        available in closed form through $A_4$ (\S\ref{sec:sub}); deriving the
        complete expansion would explain the effective constants observed at
        moderate $n$.
  \item \textbf{Order-4 construction:} find explicit closed-form solutions
        yielding $D_n=\Theta(n^{-5})$, likely requiring at least a cubic base or a
        richer exponent structure. A natural next ansatz is
        $x_n=(1+a_1/n+a_2/n^2+a_3/n^3)^{cn+\gamma}$, which introduces additional
        degrees of freedom.
  \item \textbf{Other constants (including $\gamma$):} extend the framework to
        other constants using appropriate base families. A concrete target is the
        Euler--Mascheroni constant $\gamma = \lim_{n\to\infty}(H_n-\log n)$,
        $H_n=\sum_{k\le n}1/k$; one may try to match coefficients in
        $y_n\approx\gamma + \tfrac{1}{2n} - \tfrac{1}{12n^2} + \tfrac{1}{120n^4}+\cdots$
        to enforce cancellations in $\Delta y_n$; for accurate asymptotic formulas
        for the special functions underlying such constant constructions, see
        \cite{yangtian,yangtianarxiv}.
  \item \textbf{Non-power forms:} investigate telescoping sequences outside the
        power-form ansatz.
  \item \textbf{Optimality:} characterize the maximal telescoping order achievable
        under bounded algebraic complexity.
\end{enumerate}

\section{Conclusion}
We presented an explicit closed-form two-branch construction of sequences
converging to $e^{c}$ with telescoping order $3$ (i.e.\ $|D_n|=\Theta(n^{-4})$),
obtained by enforcing coefficient-cancellation constraints within a
quadratic-base, affine-exponent family, together with explicit signed leading
constants $A_3^{(\cdot)}$ and the implied $\Theta(N^{-3})$ error rate. The broader
point is that explicit telescoping families provide clean examples where
difference-asymptotic control translates directly into error rates, making them
useful both for analysis and for didactic illustration of convergence
acceleration ideas~\cite{dlmf39}.

\subsection*{Note on this revision}
This is revision V5. Relative to V4 it corrects: the $A_3$ coefficient
formula~\eqref{eq:A3}; the constants $A_3^{(\cdot)}$ and the multiplier
$3e^{c}$ in~\eqref{eq:mainasy}; the classical baseline
(Proposition~\ref{prop:order1}, Remark~\ref{rem:correction}); the design-principle
indexing; the hypothesis of Lemma~\ref{lem:logdiff} and the signed-asymptotic
caveat (Remark~\ref{rem:theta}); Table~\ref{tab:compare}'s middle row; the
$c=0$ exception in the abstract; and several bibliographic entries. All numerical
tables were regenerated from the committed driver
(Appendix~\ref{app:code}); the one mis-rounded entry in the earlier Table~1
($8.41\times10^{-17}\to8.40\times10^{-17}$) is corrected.

\section*{Acknowledgments}
The author acknowledges the use of AI-assisted tools as editorial and programming
support during the preparation of this manuscript. All mathematical derivations,
proofs, asymptotic expansions, and numerical computations reported in the paper
were verified by the author against an independent verification driver shipped
with the paper.

\section*{Declaration of Competing Interest}
The author declares no known competing financial interests or personal
relationships that could have appeared to influence the work reported in this
paper.

\section*{Funding}
The author reports no funding for this research.

\begin{thebibliography}{99}
\bibitem{brent} R.~P. Brent, \emph{Fast multiple-precision evaluation of elementary
  functions}, J.\ ACM \textbf{23}(2) (1976), 242--251.
\bibitem{glaisher} J.~W.~L. Glaisher, \emph{On the history of Euler's constant},
  Messenger of Mathematics \textbf{1} (1872), 25--30.
\bibitem{adams} J.~C. Adams, \emph{Note on the value of Euler's constant; likewise
  on the values of the Napierian logarithms of 2, 3, 5, 7, and 10}, Proc.\ Royal
  Soc.\ London \textbf{27} (1878), 88--94.
\bibitem{mitchell} U.~G. Mitchell and M.~Strain, \emph{The number $e$}, Osiris
  \textbf{1} (1936), 476--496.
\bibitem{flajolet} R.~Sedgewick and P.~Flajolet, \emph{An Introduction to the
  Analysis of Algorithms}, booksite Chapter~4, \emph{Asymptotic Approximations},
  \url{https://aofa.cs.princeton.edu/40asymptotic/}.
\bibitem{dlmf39} NIST Digital Library of Mathematical Functions, \S3.9
  \emph{Acceleration of Convergence}, \url{https://dlmf.nist.gov/3.9}.
\bibitem{corless2023} R.~M. Corless, \emph{Devilish tricks for sequence
  acceleration}, Maple Transactions \textbf{3}(1) (2023), DOI
  \url{https://doi.org/10.5206/mt.v3i1.14777}.
\bibitem{richardson2024} C.~J. Oates, T.~Karvonen, A.~L. Teckentrup, M.~Strocchi,
  and S.~A. Niederer, \emph{Probabilistic Richardson extrapolation}, J.\ Roy.\
  Statist.\ Soc.\ Ser.\ B \textbf{87}(2) (2025), 457--479, DOI
  \url{https://doi.org/10.1093/jrsssb/qkae098} (published online 26 December 2024).
\bibitem{richardson2024arxiv} C.~J. Oates, T.~Karvonen, A.~L. Teckentrup, M.~Strocchi,
  and S.~A. Niederer, \emph{Probabilistic Richardson extrapolation},
  arXiv:2401.07562 (2024), \url{https://arxiv.org/abs/2401.07562}.
\bibitem{yangtian} Z.-H. Yang and J.-F. Tian, \emph{An accurate approximation
  formula for gamma function}, J.\ Inequalities \& Applications \textbf{2018},
  Article 56, DOI \url{https://doi.org/10.1186/s13660-018-1646-6}.
\bibitem{yangtianarxiv} Z.-H. Yang and J.-F. Tian, \emph{An accurate approximation
  formula for gamma function}, arXiv:1712.08051 (2017),
  \url{https://arxiv.org/abs/1712.08051}.
\bibitem{balditerative} J.~Bald, \emph{Telescoping representations of classical
  iterative methods and constants}, Zenodo (2025), version DOI
  \url{https://doi.org/10.5281/zenodo.17605159}; concept DOI
  \url{https://doi.org/10.5281/zenodo.17605158}.
\end{thebibliography}

\appendix
\section{Reproducibility Code (Python)}\label{app:code}
This appendix provides a self-contained script reproducing the high-precision
numerical validation in \S5: computation of $x_n^{(\mathrm{I})}$,
$x_n^{(\mathrm{II})}$, differences $D_n=x_{n+1}-x_n$, log--log slope estimates of
$|D_n|$ versus $n$, and sample convergence values for $c=1,2,-1$.

\emph{Dependencies.} \texttt{mpmath} (arbitrary precision). The full verification
driver (\texttt{code/verify\_hoc.py}), which additionally checks the symbolic
coefficients~\eqref{eq:A3}--\eqref{eq:A4} and \emph{asserts every claimed sign and
constant}, is available in the reproduction package.

\begin{lstlisting}[basicstyle=\ttfamily\footnotesize,breaklines=true,frame=single]
import mpmath as mp

def xI(n, c):
    n = mp.mpf(n); s3 = mp.sqrt(3)
    return mp.power(1 + 1/n + s3/(6*n**2), c*n + c*(mp.mpf(1)/2 - s3/6))

def xII(n, c):
    n = mp.mpf(n); s3 = mp.sqrt(3)
    return mp.power(1 + 1/n - s3/(6*n**2), c*n + c*(mp.mpf(1)/2 + s3/6))

def slope(f, c, n0=100, n1=10000, step=25):
    ns = list(range(n0, n1+1, step))
    xs = [mp.log(n) for n in ns]
    ys = [mp.log(mp.fabs(f(n+1, c) - f(n, c))) for n in ns]
    m = len(ns); mx = mp.fsum(xs)/m; my = mp.fsum(ys)/m
    cov = mp.fsum([(xs[i]-mx)*(ys[i]-my) for i in range(m)])
    var = mp.fsum([(xs[i]-mx)**2 for i in range(m)])
    return cov/var, my - (cov/var)*mx

if __name__ == "__main__":
    mp.mp.dps = 150
    c = mp.mpf(1); target = mp.e**c
    print("slope(I) :", slope(xI, c))
    print("slope(II):", slope(xII, c))
    for n in [100, 500, 1000, 5000]:
        print(n, mp.fabs(xI(n,c)-target), mp.fabs(xI(n+1,c)-xI(n,c)),
                 mp.fabs(xII(n,c)-target), mp.fabs(xII(n+1,c)-xII(n,c)))
    # leading-constant check:  n^4 * D_n  ->  -3 e^c A3
    for f, A3 in [(xI, (2*mp.sqrt(3)-3)/72), (xII, -(3+2*mp.sqrt(3))/72)]:
        n = 10**5
        print("n^4*D_n ->", (f(n+1,c)-f(n,c))*n**4, " predicted", -3*mp.e*A3)
\end{lstlisting}

\section{Verification Summary}\label{app:verify}
The companion driver (\texttt{code/verify\_hoc.py}) regenerates every numeral in
this paper and asserts each claimed sign and constant (and exits with nonzero
status on any failed assertion); it reports \texttt{ALL ASSERTIONS PASS}.
Symbolically it confirms \eqref{eq:c0}--\eqref{eq:A4} (including the general
$A_4$ identity), $A_1=A_2=0$ on both branches, and the constants
$A_3^{(\mathrm{I})}=c(2\sqrt3-3)/72$, $A_3^{(\mathrm{II})}=-c(3+2\sqrt3)/72$,
$A_4^{(\mathrm{I},\mathrm{II})}=c(39\mp25\sqrt3)/720$. Numerically it confirms
$\lim x_n=e^{c}$, $n^4D_n\to-3e^{c}A_3$ and $n^3(x_n-e^{c})\to e^{c}A_3$ with the
correct signs, every tabulated value in Tables~\ref{tab:c1}--\ref{tab:gen}, the
regression slopes and $R^2$, and the corrected classical baseline
$D_n\sim e^{ab}a^2b/(2n^2)$. The exact algebraic core (the branch cancellations
and the constants $A_3$, $A_4$) is machine-checked in Lean~4
(\texttt{lean/ExpHOC}); the gate reports \texttt{PASS} (13 theorems, standard
axioms only).

The complete reproduction package --- manuscript source, driver and its saved
output, the Lean development, and the independent audit reports --- is available
at \url{https://github.com/jpbald93/telescoping-exponential-approximants}.

\end{document}
